\documentclass[sigconf,nonacm]{acmart}

\usepackage{booktabs}
\usepackage{amsmath}
\usepackage{graphicx}
\usepackage{listings}
\usepackage{xcolor}
\usepackage{pifont}
\newcommand{\cmark}{\textcolor{green!55!black}{\ding{51}}}
\newcommand{\xmark}{\textcolor{red!70!black}{\ding{55}}}

% Allow more flexible line breaking to avoid overfull \hbox in two-column.
\sloppy
\emergencystretch=3em

\title[Learning Optimization Skills for GPU Kernels]{Learning Optimization Skills for GPU Kernels: A Knowledge--Environment--Objective Framework}

\author{Haibin Lai}
\affiliation{%
  \institution{Southern University of Science and Technology}
  \city{Shenzhen}
  \country{China}
}
\email{12211612@mail.sustech.edu.cn}

\begin{document}

\begin{abstract}
Optimizing GPU kernels for high-performance computing (HPC) workloads remains a labor-intensive process that relies on implicit expert knowledge. We discover that this knowledge has a \emph{two-level structure}: a compact set of optimization \emph{strategies} (e.g., tiled GEMM with online reduction, temporal blocking, SMEM-cached streaming) that transfer across operators, and a longer tail of operator-specific \emph{implementations} (e.g., correct butterfly indexing for FFT) that must be supplied as verified code. We formalize this through a Knowledge--Environment--Objective (KEO) framework where an LLM-driven agent interacts with hardware profilers (NVIDIA Nsight Compute), domain knowledge, and multi-scale benchmarks. Experiments on 10~operators covering all 7 Berkeley Dwarfs on an NVIDIA A100 GPU show that 5 strategies cover 8~operators achieving 1.2--6.0$\times$ speedup, that strategy-level transfer succeeds across operators (3/3) while implementation-level transfer fails (1/1), and that NCU-guided diagnosis turns a stuck 0.07$\times$ kernel into a 1.44$\times$ success in one iteration. A complete end-to-end case study on Flash-GMM achieves 1.56$\times$ speedup and 32$\times$ VRAM reduction, with $\sim\!10\times$ NCU-guided kernel refinement reusing existing skills.
\end{abstract}

\keywords{GPU optimization, kernel compilation, optimization skills, LLM agents, hardware profiling}

\maketitle

% ============================================================================
\section{Introduction}
% ============================================================================

GPU kernel optimization is one of the most expertise-intensive tasks in high-performance computing. Achieving near-peak performance on modern GPUs requires navigating a complex space of decisions: memory hierarchy exploitation, parallelism patterns, hardware-specific features (Tensor Cores, shared memory), and algorithm-level transformations (online algorithms, kernel fusion). This knowledge is typically \emph{implicit}---residing in the minds of expert programmers---and \emph{fragmented}---scattered across papers, library implementations, and tribal knowledge.

Recent work has explored two approaches to automate this process. \emph{Rule-based systems} encode expert knowledge as static heuristics or search templates, but cannot generalize beyond their predefined rules. \emph{Search-based systems} explore the optimization space through autotuning or reinforcement learning, but treat each kernel independently, discarding the structural similarity between optimization strategies.

We make a key observation: \textbf{HPC kernel optimization has reusable structure}. The same optimization \emph{pattern}---such as ``tile over the reduction dimension, compute partial results via GEMM in shared memory, apply an online algorithm to avoid materializing intermediates''---applies to KMeans (distance + argmin), N-Body (force + accumulation), and GMM (likelihood + logsumexp), despite these being entirely different computational problems. We call such reusable patterns \emph{optimization skills}.

This observation leads to three refined claims:

\begin{enumerate}
    \item \textbf{Two-Level Skill Discovery.} Optimization knowledge
    decomposes into a \emph{strategy} layer (transferable, learnable
    from trajectories: e.g.\ ``SMEM fusion of multiple passes'',
    ``higher radix'') and an \emph{implementation} layer
    (operator-specific, must be supplied as verified code
    templates: e.g.\ correct butterfly indexing).
    LLM agents can reliably discover the former but consistently fail
    at the latter.
    \item \textbf{Skill Transfer.} Strategies discovered on one
    kernel transfer to unseen kernels that share computational
    structure, even though their implementations do not.
    \item \textbf{Skill Compactness.} A small number of strategies
    (5--10) cover a large portion of the optimization space across
    the 7 Berkeley Dwarfs, multiplied by a longer tail of
    operator-specific implementations.
\end{enumerate}

To validate these claims, we build an LLM-driven agent that interacts with three interfaces:
\begin{itemize}
    \item \textbf{Knowledge} (K): retrieval, summarization, and accumulation of optimization patterns from papers, libraries, and prior experience.
    \item \textbf{Environment} (E): hardware profiling via NVIDIA Nsight Compute (NCU), providing bandwidth utilization, compute throughput, Tensor Core usage, and occupancy.
    \item \textbf{Objective} (O): multi-scale benchmarking with version tracking, regression detection, and library ceiling comparison.
\end{itemize}

We evaluate on 10 operators spanning all 7 Berkeley Dwarfs (dense/sparse linear algebra, spectral methods, N-body, structured/unstructured grids, MapReduce) on an NVIDIA A100 GPU. Our results show:

\begin{itemize}
    \item 5 strategies cover 10 operators across 7 Dwarfs, achieving 1.18--6.04$\times$ speedup where applicable.
    \item The \texttt{tiled\_gemm\_online\_reduce} strategy, discovered on KMeans, transfers directly to N-Body (2.19$\times$) and GMM (1.44$\times$).
    \item NCU-guided diagnosis enables the agent to progress from 0.07$\times$ (broken kernel) to 1.44$\times$ (correct kernel) in a single iteration, compared to blind trial-and-error.
    \item The agent correctly identifies 4 out of 5 cases where IO optimization is \emph{not} effective (Softmax, Conv2D, SpMV, KMeans-scalar), avoiding wasted effort.
    \item Our FFT case study (\S\ref{sec:fft-case}) shows that across
    9 kernel rewrites the LLM correctly chose every strategy yet
    never produced correct butterfly indexing, motivating the
    two-level skill structure.
\end{itemize}

% ============================================================================
\section{Background and Motivation}
% ============================================================================

\subsection{GPU Memory Hierarchy and the Roofline Model}
\label{sec:bg-roofline}

A modern data-center GPU exposes a steep memory hierarchy: registers
($\sim\!64\,$KB per SM, $\sim\!1\,$cycle), shared memory / L1
($\sim\!192\,$KB per SM, $\sim\!20\,$cycle), L2 cache (40--50\,MB
chip-wide, $\sim\!200\,$cycle), and HBM
(40--80\,GB, 2--5\,TB/s, $\sim\!500\,$cycle). The cost of moving a
byte across this hierarchy varies by three orders of magnitude, so
\emph{where} an intermediate tensor lives often matters more than how
many FLOPs are spent on it.

The roofline model formalizes this trade-off. For a kernel with
arithmetic intensity $I = \mathrm{FLOPs}/\mathrm{HBM\;bytes}$ on a
device with peak compute $\pi$ and peak bandwidth $\beta$, achievable
performance is bounded by $\min(\pi,\,I\beta)$, with the
\emph{balance point} $I^\star = \pi/\beta$ separating
\emph{memory-bound} from \emph{compute-bound} regimes. Because peak
TFLOPS has grown faster than HBM bandwidth across hardware
generations---$I^\star \approx 10$ on A100, $\approx 14$ on H200---an
ever larger fraction of HPC kernels falls below the balance point and
becomes memory-bound. For these kernels, reducing HBM traffic is the
dominant lever.

\subsection{IO-Aware Kernel Design}

FlashAttention~\cite{dao2022flashattention} popularized the
\emph{IO-aware} discipline of fusing memory-bound passes via online
algorithms: never materialize an $N{\times}N$ attention score matrix
to HBM; stream it through SMEM and update an online softmax in
registers. Subsequent work---Flash-KMeans~\cite{flashkmeans2026},
Flash-GMM (this paper, \S\ref{sec:flash-gmm})---confirms the recipe:
identify a materialized intermediate, replace it with an online
algorithm that maintains $O(1)$ register state per row, and tile the
input/parameters so the inner data stays on chip. This recipe trades
recomputation FLOPs for HBM bytes, which is profitable whenever the
kernel sits below the roofline balance point.

\subsection{The Berkeley Dwarfs as a Coverage Frame}
\label{sec:bg-dwarfs}

To avoid drawing conclusions from a single workload family, we
benchmark across the seven \emph{Berkeley Dwarfs}~\cite{asanovic2009dwarfs}---an
established taxonomy of HPC computational patterns: dense linear
algebra, sparse linear algebra, spectral methods, N-body, structured
grids, unstructured grids, and MapReduce. Each dwarf stresses
different combinations of memory access patterns, parallelism, and
on-chip resource pressure, so a framework that succeeds across all
seven is unlikely to be over-fit to attention-like workloads.

\subsection{The IO Optimization Hypothesis}

Building on the roofline view, the natural hypothesis is that any
memory-bound kernel benefits from IO-aware fusion. We tested this on
all seven dwarfs and found that \textbf{IO reduction is necessary but
not sufficient}:

\begin{table}[h]
\centering
\caption{IO reduction vs.\ actual speedup on A100. IO optimization fails when the operator is compute-bound, data fits in L2 cache, or the optimization destroys GEMM parallelism.}
\label{tab:io-vs-speedup}
\setlength{\tabcolsep}{3pt}
\small
\begin{tabular}{lccl}
\toprule
Operator & IO $\downarrow$ & Speedup & Root Cause of Mismatch \\
\midrule
CrossEntropy    & 80\%  & \textbf{6.04$\times$} & --- mem-bound, data $\gg$ L2 \\
LayerNorm       & 50\%  & \textbf{4.58$\times$} & --- memory-bound \\
Stencil2D       & 99\%  & \textbf{4.74$\times$} & --- temporal tiling \\
N-Body          & 100\% & \textbf{2.19$\times$} & --- SMEM tiling kills $N^2$ \\
KMeans (tl.dot) & 94\%  & \textbf{1.44$\times$} & --- GEMM preserved \\
FFT             & 93\%  & \textbf{1.32$\times$} & --- radix-2 butterfly \\
Graph Laplacian & 17\%  & \textbf{1.18$\times$} & --- fused scatter-gather \\
\midrule
KMeans (scalar) & 94\% & 0.07$\times$ & Lost Tensor Core \\
Softmax         & 65\% & 0.86$\times$ & Data fits L2 \\
Conv2D          & 86\% & 0.57$\times$ & cuDNN already optimal \\
SpMV            & 0\%  & 0.76$\times$ & Wrong sparsity pattern \\
\bottomrule
\end{tabular}
\end{table}

This motivates a multi-dimensional analysis beyond IO complexity alone.

\subsection{The Skill Hypothesis}

We observe that successful optimizations share common \emph{skills} across different operators:

\begin{table}[h]
\centering
\caption{5 optimization skills cover 10 operators across 7 Berkeley Dwarfs.}
\label{tab:skills}
\setlength{\tabcolsep}{3.5pt}
\small
\begin{tabular}{lll}
\toprule
Skill & Operators & Dwarfs \\
\midrule
Online streaming reduce  & CE, LN, GMM     & MapReduce \\
Tiled GEMM + online red. & KMeans, N-Body  & Dense LA, N-Body \\
Temporal tiling          & Stencil2D       & Structured Grid \\
SMEM butterfly fusion    & FFT             & Spectral \\
Fused scatter-gather     & Graph Laplacian & Unstructured Grid \\
\bottomrule
\end{tabular}
\end{table}

This compactness suggests that learning a small skill library could dramatically reduce the search space for new kernels.

% ============================================================================
\section{The KEO Framework}
\label{sec:keo}
% ============================================================================

\subsection{Formalization}

We formalize kernel optimization as an iterative process over three
interfaces (Figure~\ref{fig:keo}):

\begin{figure}[t]
\centering
\includegraphics[width=0.92\linewidth]{figures/keo_architecture.pdf}
\caption{KEO architecture: the LLM agent interacts with Knowledge
(patterns, library analysis), Environment (NCU profiling), and
Objective (benchmarking, version tracking) interfaces. Knowledge
stores two-level skills: transferable strategies and verified
implementations.}
\label{fig:keo}
\end{figure}

\textbf{Knowledge ($\mathcal{K}$).} A growing library of optimization skills $\mathcal{S} = \{s_1, \ldots, s_k\}$ where each skill $s = (c, p, t)$ consists of:
\begin{itemize}
    \item \emph{Condition} $c$: when the skill applies (e.g., ``operator is memory-bound with materialized $N \times K$ intermediate'')
    \item \emph{Pattern} $p$: the optimization structure (e.g., ``tile over $K$, compute partial GEMM via \texttt{tl.dot}, apply online argmin'')
    \item \emph{Transformation} $t$: the code transformation (e.g., a Triton kernel template)
\end{itemize}

\textbf{Environment ($\mathcal{E}$).} Hardware profiling metrics $e = (\text{BW}_\%, \text{Compute}_\%, \text{TC}, \text{Occ}_\%, \ldots)$ obtained from NVIDIA Nsight Compute (NCU), providing ground-truth utilization data.

\textbf{Objective ($\mathcal{O}$).} Performance evaluation $o = (\text{speedup}, \text{correctness}, \text{version\_history})$ with multi-scale benchmarking and library ceiling comparison.

The agent loop is:
\begin{equation}
    a_t = \pi(s_{1:t-1}, e_{1:t-1}, o_{1:t-1}; \mathcal{K})
\end{equation}
where $\pi$ is the LLM policy, $a_t$ is the action at step $t$ (analyze, fuse, generate kernel, profile, etc.), and $\mathcal{K}$ is updated after successful optimizations.

\subsection{Skill Lifecycle}

\begin{enumerate}
    \item \textbf{Discovery.} From an optimization trajectory $(a_1, e_1, o_1, \ldots, a_T, e_T, o_T)$ that achieves speedup $> 1\times$, extract the key (condition, pattern, transformation) triple.
    \item \textbf{Abstraction.} Generalize the skill by replacing operator-specific details with structural features (e.g., ``$N \times K$ intermediate'' $\to$ ``materialized pairwise computation'').
    \item \textbf{Transfer.} Match new operators against skill conditions; apply the highest-scoring skill.
    \item \textbf{Evolution.} Update skill success statistics; refine conditions based on failure cases (e.g., ``only when data $>$ L2 cache'').
\end{enumerate}

% ============================================================================
\section{System Design}
% ============================================================================

\subsection{Agent Architecture}

The agent exposes 20 tools organized by KEO interface:

\begin{itemize}
    \item \textbf{K tools:} \texttt{retrieve\_pattern}, \texttt{pre\_check}
          (7-dim static analysis), \texttt{show\_actions}.
    \item \textbf{E tools:} \texttt{ncu\_profile} (NCU counters
          + torch.profiler), \texttt{library\_ceiling},
          \texttt{compare\_profile}, \texttt{occupancy\_analysis}.
    \item \textbf{O tools:} \texttt{benchmark\_kernel} (version
          tracking + rollback), \texttt{compile\_and\_test},
          \texttt{autotune\_kernel}, \texttt{verify}.
    \item \textbf{Graph tools:} \texttt{analyze}, \texttt{fuse\_ops},
          \texttt{fuse\_and\_online}, \texttt{generate\_kernel}
          (Triton/CUDA codegen).
\end{itemize}

\subsection{NCU-Driven Diagnosis}

Hardware profiling provides ground-truth signals that static analysis cannot:

\begin{table}[h]
\centering
\caption{NCU metrics enable precise diagnosis of optimization failures.}
\label{tab:ncu}
\setlength{\tabcolsep}{3.5pt}
\small
\begin{tabular}{lcccl}
\toprule
Operator & DRAM\% & SM\% & Warps\% & Diagnosis \\
\midrule
CrossEntropy   & 77 & 68 & 66 & Memory-bound \\
LayerNorm      & 80 & 47 & 95 & Memory-bound \\
KMeans (scalar)& 0  & 75 & 99 & Scalar loop, no DRAM \\
KMeans (tl.dot)& 1  & 15 & 35 & Compute-bound, TC on \\
Softmax        & 7  & 49 & 82 & Data fits L2 \\
\bottomrule
\end{tabular}
\end{table}

The signal \texttt{warps=99\%, DRAM=0\%} for KMeans-scalar is \emph{impossible to detect without runtime profiling}---it reveals that the SM is 99\% busy doing scalar arithmetic (not memory accesses), despite the kernel having 94\% IO reduction on paper.

\subsection{Version Tracking and Regression Detection}

LLM-generated kernels can regress (e.g., 5.83$\times$ $\to$ 3.02$\times$ on a rewrite attempt). We maintain:
\begin{itemize}
    \item A version history with speedup per version
    \item Automatic rollback when new speedup $< 0.95 \times$ best
    \item Convergence detection: \texttt{done} only after \texttt{ncu\_profile} + \texttt{library\_ceiling} + \texttt{autotune}
\end{itemize}

% ============================================================================
\section{Evaluation}
% ============================================================================

\subsection{Experimental Setup}

\textbf{Hardware.} NVIDIA A100 80GB PCIe (2039 GB/s HBM, 19.5 TFLOPS FP32, 312 TFLOPS BF16 TC).

\textbf{Operators.} 10 operators covering all 7 Berkeley Dwarfs: CrossEntropy, LayerNorm, Softmax, KMeans, FFT, Conv2D, Stencil2D, N-Body, SpMV, Graph Laplacian.

\textbf{Agent.} GPT-5.2 via Copilot proxy, max 25--30 steps per operator.

\textbf{Baselines.} (1) Materialized PyTorch baseline (explicit intermediates), (2) Library ceiling (cuDNN/cuFFT/cuBLAS/ATen), (3) Static IO analysis only (no profiling).

\subsection{Skill Transfer (Claim 2)}

We separate \emph{strategy transfer} (apply the abstract pattern,
let the agent re-derive the implementation for the new operator)
from \emph{implementation transfer} (copy the source kernel's code
verbatim, retargeted only by name substitution).
Table~\ref{tab:transfer} reports both modes on three target
operators.

\begin{table}[h]
\centering
\caption{Cross-operator skill transfer. Strategy transfer (3/3
positive) vs.\ verbatim implementation transfer (1/1 regression):
KMeans's \texttt{tl.dot} Triton kernel, copied to N-Body, regresses
to 0.68$\times$ because $d{=}3$ does not fit a GEMM tile.}
\label{tab:transfer}
\small
\begin{tabular}{llccc}
\toprule
Source & Target & w/o Skill & Strategy & Impl.\ copy \\
\midrule
KMeans       & N-Body     & 0.07$\times$ & \textbf{2.19$\times$} & 0.68$\times$ \\
KMeans       & GMM E-step & 0.11$\times$ & \textbf{1.44$\times$} & --- \\
CrossEntropy & LayerNorm  & ---          & \textbf{4.87$\times$} & --- \\
\bottomrule
\end{tabular}
\end{table}

\textbf{Finding.} Strategy-layer transfer succeeds in all three
cases, while a direct implementation copy regresses N-Body to
0.68$\times$: KMeans's \texttt{tl.dot} tile assumes a high inner
dimension ($d{\geq}32$), but N-Body has $d{=}3$, so the GEMM tile
is mostly padding. This mirrors the FFT observation
(\S\ref{sec:fft-case})---transferable knowledge lives at the
strategy layer; implementations must be re-derived per operator.

\subsection{Skill Coverage (Claim 3)}

\begin{table}[h]
\centering
\caption{5 skills cover 10 operators across 7 Dwarfs.}
\label{tab:coverage}
\setlength{\tabcolsep}{4pt}
\small
\begin{tabular}{lcccc}
\toprule
\# Skills & Operators Covered & Dwarfs Covered & Avg.\ Speedup \\
\midrule
1 & 3 (CE, LN, GMM) & 1 & 5.0$\times$ \\
2 & 5 (+KMeans, N-Body) & 3 & 3.8$\times$ \\
3 & 6 (+Stencil) & 4 & 3.9$\times$ \\
4 & 7 (+FFT) & 5 & 3.5$\times$ \\
5 & 8 (+Graph Lap.) & 6 & 3.2$\times$ \\
\bottomrule
\end{tabular}
\end{table}

\subsection{NCU Impact (Profiling-Guided vs.\ Blind)}

Directly comparing agent outcomes with and without NCU access
shows that hardware profiling is what turns a stuck trajectory
into a converging one.

\begin{table}[h]
\centering
\caption{Impact of NCU-driven diagnosis on optimization outcome.}
\label{tab:ncu-impact}
\setlength{\tabcolsep}{3.5pt}
\small
\begin{tabular}{lll}
\toprule
Operator & Blind (no NCU) & NCU-Guided \\
\midrule
KMeans       & 0.07$\times$ stuck   & \textbf{1.44$\times$}, 1 iter. \\
Softmax      & 0.50$\times$ wasted  & 0.86$\times$, flagged unfixable \\
CrossEntropy & 5.83$\times$, 3 iter.& 5.83$\times$, confirmed in 1 iter. \\
\bottomrule
\end{tabular}
\end{table}



\subsection{Ablation: Tool Configuration Impact}

To isolate each tool's contribution, we run four configurations
on four operators (Table~\ref{tab:ablation}).
For memory-bound operators (CrossEntropy, LayerNorm), the template
alone suffices and additional tools provide only confirmation.
For KMeans, which requires an algorithmic fix (scalar loop $\to$
\texttt{tl.dot} GEMM), the full toolkit achieves a 20$\times$
improvement over the template: NCU diagnoses the problem,
\texttt{retrieve\_pattern} provides the fix, and
\texttt{analyze\_library} enables a proactive strategy that skips
the broken template entirely.

\begin{table}[h]
\centering
\caption{Ablation: speedup under four tool configurations.
Config~D lets the agent skip the broken template on KMeans.}
\label{tab:ablation}
\setlength{\tabcolsep}{3.5pt}
\small
\begin{tabular}{lcccc}
\toprule
Operator & A: Template & B: +NCU & C: +Ceiling & D: Full \\
\midrule
KMeans       & 0.07$\times$ & 0.07$\times$ & 0.07$\times$ & \textbf{1.43$\times$} \\
CrossEntropy & 5.84$\times$ & 5.89$\times$ & 5.84$\times$ & 5.83$\times$ \\
LayerNorm    & 4.83$\times$ & 4.85$\times$ & 4.84$\times$ & 4.85$\times$ \\
Softmax      & 1.02$\times$ & 1.01$\times$ & 1.01$\times$ & 1.02$\times$ \\
\bottomrule
\end{tabular}
\end{table}

\begin{figure}[h]
\centering
\includegraphics[width=0.92\linewidth]{figures/ablation.pdf}
\Description{Bar chart comparing speedups under four agent tool configurations across four operators; KMeans only succeeds under the full configuration, while CrossEntropy, LayerNorm, and Softmax are insensitive to the configuration.}
\caption{Ablation study: the full KEO toolkit is decisive only for
operators that require algorithmic repair (KMeans); memory-bound
operators succeed with the template alone.}
\label{fig:ablation}
\end{figure}

% KEO architecture figure moved to §3 (Figure 1)
\end{figure}


% ============================================================================
\section{Case Study: Flash-GMM}
\label{sec:flash-gmm}
% ============================================================================

To validate the KEO framework on a complete end-to-end application
(rather than isolated kernels), we apply it to derive \emph{Flash-GMM},
an IO-aware Gaussian Mixture Model EM solver that extends FlashAttention's
online-algorithm style to the EM update. The agent autonomously
discovered the two-pass fusion structure, the SMEM tiling strategy, and
ten distinct CUDA-level micro-optimizations distinguishing v1 from v2.

\subsection{Algorithm and IO Saving}

Standard GMM EM materializes both the log-likelihood matrix
$L \in \mathbb{R}^{N \times K}$ and the responsibility matrix
$\gamma \in \mathbb{R}^{N \times K}$ to HBM, incurring $\Theta(NK)$ traffic
per iteration. Guided by the \texttt{online streaming reduce} skill, the
agent derived a two-pass formulation: (i)~\textbf{Flash log-normalizer}
streams centroid tiles through SMEM and maintains an online log-sum-exp
in registers, outputting only $\mathtt{log\_norm} \in \mathbb{R}^N$;
(ii)~\textbf{Flash accumulate stats} recomputes $\gamma$ on-chip and
accumulates sufficient statistics $(n_k, s_k, sq_k)$ directly. This
eliminates \emph{all} $\Theta(NK)$ traffic. The closed-form saving
$\approx 2K/d + 1$ is validated empirically:

\begin{table}[h]
\centering
\caption{Theoretical HBM traffic per EM iteration.}
\label{tab:gmm-io}
\setlength{\tabcolsep}{4pt}
\small
\begin{tabular}{rrrrrrr}
\toprule
$N$ & $K$ & $d$ & $K/d$ & Std IO & Flash IO & Saving \\
\midrule
4{,}096  & 64  & 32  & 2.0 & 5.28 MB  & 1.10 MB  & \textbf{4.80$\times$}  \\
4{,}096  & 256 & 32  & 8.0 & 17.96 MB & 1.25 MB  & \textbf{14.39$\times$} \\
8{,}192  & 512 & 64  & 8.0 & 71.83 MB & 4.99 MB  & \textbf{14.40$\times$} \\
16{,}384 & 256 & 128 & 2.0 & 84.41 MB & 17.57 MB & \textbf{4.81$\times$}  \\
\bottomrule
\end{tabular}
\end{table}

\subsection{End-to-End GPU Performance}

We benchmark on A100 80GB across 12 configurations. The
\texttt{pre\_check} tool flags $d{=}128$ as memory-bound (the $N{\times}K$
intermediates well exceed the 40\,MB L2); for $d{=}64$ the M-step's GEMM
is compute-bound and dominated by cuBLAS Tensor Cores.

\begin{table}[h]
\centering
\caption{Flash-GMM E+M vs.\ standard GMM on A100 (selected configurations).}
\label{tab:gmm-gpu}
\setlength{\tabcolsep}{3.5pt}
\small
\begin{tabular}{rrrrrrr}
\toprule
$N$ & $K$ & $d$ & Std (ms) & Flash (ms) & Speedup & VRAM $\downarrow$ \\
\midrule
65{,}536 & 64   & 128 & 5.42  & 3.58  & \textbf{1.51$\times$} & 2.0$\times$ \\
65{,}536 & 256  & 128 & 21.65 & 13.84 & \textbf{1.56$\times$} & 5.0$\times$ \\
65{,}536 & 512  & 128 & 42.54 & 27.61 & \textbf{1.54$\times$} & \textbf{8.9$\times$} \\
65{,}536 & 1024 & 64  & 45.41 & 51.67 & 0.88$\times$          & \textbf{32.0$\times$} \\
32{,}768 & 1024 & 64  & 30.01 & 30.61 & 0.98$\times$          & \textbf{31.2$\times$} \\
16{,}384 & 128  & 64  & 1.99  & 2.75  & 0.73$\times$          & 4.9$\times$ \\
\bottomrule
\end{tabular}
\end{table}

\textbf{Findings.} (1)~When \texttt{pre\_check} signals memory-bound
($d{=}128$), Flash-GMM consistently delivers $1.51$--$1.56\times$
speedup, matching the roofline prediction.
(2)~Even when wall-clock is neutral ($d{=}64$, compute-bound), VRAM
shrinks by up to $32\times$, enabling otherwise infeasible problem sizes.
(3)~Correctness: 20-iteration convergence on all 12 configurations
yields final log-likelihood difference $0.00\mathrm{e}{+}00$ vs.\ the
standard implementation.

\subsection{CPU Path: Skill Reuse on Different Hardware}

The same KEO loop, with \texttt{Environment} swapped from NCU to
\texttt{perf} counters and codegen retargeted from CUDA to
C++/OpenMP, automatically produced a CPU implementation. The
\texttt{tiled\_gemm\_online\_reduce} skill's tile sizes were
re-selected by autotune for L2 cache.

\begin{table}[h]
\centering
\caption{C++ OpenMP (8 cores) vs.\ PyTorch ATen/MKL (8 cores). Skill
transfer to CPU yields linear scaling that PyTorch fails to achieve.}
\label{tab:gmm-cpu}
\setlength{\tabcolsep}{4pt}
\small
\begin{tabular}{lrrr}
\toprule
$(N, K, d)$ & PyTorch & C++ Std & Speedup \\
\midrule
(4{,}096, 8, 32)     & 4.21 ms     & \textbf{0.75 ms}    & \textbf{5.6$\times$} \\
(4{,}096, 64, 32)    & 18.01 ms    & \textbf{5.40 ms}    & \textbf{3.3$\times$} \\
(16{,}384, 64, 64)   & 86.52 ms    & \textbf{37.97 ms}   & \textbf{2.3$\times$} \\
(16{,}384, 128, 64)  & 185.84 ms   & \textbf{75.05 ms}   & \textbf{2.5$\times$} \\
(65{,}536, 256, 128) & 2854.02 ms  & \textbf{1148.78 ms} & \textbf{2.5$\times$} \\
\bottomrule
\end{tabular}
\end{table}

Parallel scaling further confirms skill-level transfer: C++/OpenMP
reaches $6.3$--$7.9\times$ on 8 cores while PyTorch saturates at
$1.9$--$3.6\times$.

\subsection{Kernel Iteration: v1 to v2}

Within the GPU codegen path, the agent issued three
\texttt{ncu\_profile} $\to$ rewrite cycles. Hardware counters revealed:
(i)~\texttt{logf} dominating the inner loop, (ii)~non-coalesced
Mahalanobis reads, (iii)~SMEM \texttt{atomicAdd} contention in the
M-step. Each diagnosis triggered a targeted transformation
(precomputed \texttt{inv\_var}/\texttt{log\_coeff} hoisted out of the
hot loop, lane-striped warp-collaborative Mahalanobis distance,
per-warp accumulators) drawn from the existing skill library.

\begin{table}[h]
\centering
\caption{Kernel-internal speedup from v1 to v2, attributable to
NCU-guided diagnosis on the \emph{same} algorithmic structure.}
\label{tab:gmm-v1v2}
\setlength{\tabcolsep}{4pt}
\small
\begin{tabular}{rrrrrr}
\toprule
$N$ & $K$ & $d$ & v1 vs.\ Std & v2 vs.\ Std & Kernel gain \\
\midrule
65{,}536 & 64  & 128 & 0.14$\times$ & \textbf{1.51$\times$} & \textbf{10.8$\times$} \\
65{,}536 & 256 & 128 & 0.16$\times$ & \textbf{1.56$\times$} & \textbf{9.8$\times$}  \\
65{,}536 & 512 & 128 & 0.15$\times$ & \textbf{1.54$\times$} & \textbf{10.3$\times$} \\
\bottomrule
\end{tabular}
\end{table}

\textbf{No new skill was authored to drive the v1$\to$v2 jump}---a
direct demonstration of Claim~3 (skill compactness).

\begin{figure}[h]
\centering
\includegraphics[width=0.92\linewidth]{figures/gmm_gpu_speedup.pdf}
\Description{Scatter plot of Flash-GMM speedup versus standard GMM across configurations on A100; configurations with d=128 form an upper cluster around 1.5x speedup, while d=64 configurations cluster near 1x.}
\caption{Flash-GMM speedup vs.\ standard GMM on A100. The agent's
\texttt{pre\_check} tool predicts memory-bound regimes ($d{=}128$,
shown as the upper cluster) where IO optimization is profitable.}
\label{fig:gmm-gpu}
\end{figure}

\begin{figure}[h]
\centering
\includegraphics[width=0.92\linewidth]{figures/skills_pareto.pdf}
\Description{Pareto curve showing the cumulative number of operators and Berkeley Dwarfs covered as optimization skills are added; five skills cover eight of ten operators across six of seven dwarfs.}
\caption{Skill-coverage Pareto curve: cumulative \# operators
covered as skills are added. Five skills suffice to cover 8 of 10
operators across 6 of 7 Berkeley Dwarfs.}
\label{fig:skills-pareto}
\end{figure}

% ============================================================================
\section{Case Study: FFT and the Two-Level Skill Structure}
\label{sec:fft-case}
% ============================================================================

While Flash-GMM demonstrates a successful end-to-end skill
application, our FFT experiment is the most informative
\emph{negative} result: the agent reached a 1.32$\times$
speedup (Tab.~\ref{tab:io-vs-speedup}) but only after the
human-supplied implementation template was retrieved; nine
preceding LLM rewrites were all numerically incorrect.
This contrast separates two layers of optimization knowledge that
prior work conflates.

\subsection{What the Agent Got Right}

Across the 9 rewrites, the LLM consistently and correctly
selected the high-level strategy:
\emph{(i)}~switch from radix-2 to radix-4 to halve the number of
passes;
\emph{(ii)}~fuse 12 successive butterfly stages in SMEM rather
than 10, eliminating $\sim\!50\,$MB of HBM round-trips;
\emph{(iii)}~retarget from Triton to CUDA after
\texttt{retrieve\_pattern} indicated that warp-shuffle butterfly
permutations cannot be expressed in Triton.
Each decision matched the canonical cuFFT design. The strategic
reasoning is reproducible across cold-start agent sessions and is
indistinguishable from expert reasoning.

\subsection{What the Agent Got Wrong}

\begin{table}[h]
\centering
\caption{Strategy-level vs.\ implementation-level skills observed
across 9 FFT kernel rewrites. The LLM scores 4/4 on strategy and
0/3 on implementation.}
\label{tab:fft-2level}
\setlength{\tabcolsep}{4pt}
\small
\begin{tabular}{llc}
\toprule
Layer & Decision/Detail & Correct? \\
\midrule
Strategy       & radix-4 vs.\ radix-2                & \cmark \\
Strategy       & Cooley--Tukey mathematics           & \cmark \\
Strategy       & Triton vs.\ CUDA target             & \cmark \\
Strategy       & Add radix-2 prelude after debug log & \cmark \\
\midrule
Implementation & Twiddle index uses global $N$, not tile $N$ & \xmark \\
Implementation & Cooley--Tukey in-place vs.\ Stockham        & \xmark \\
Implementation & Bit-reverse permutation formula             & \xmark \\
\bottomrule
\end{tabular}
\end{table}

\begin{figure}[h]
\centering
\includegraphics[width=0.92\linewidth]{figures/fft_two_level.pdf}
\Description{Two-row matrix showing strategy decisions (top, all green checkmarks) versus implementation details (bottom, all red crosses) across 9 FFT rewrites; right side shows aggregate score 4/4 strategy, 0/3 implementation.}
\caption{FFT two-level skill structure across 9 LLM rewrites:
the agent scores 4/4 on strategy decisions but 0/3 on
implementation details. This motivates the strategy/implementation
decomposition.}
\label{fig:fft-2level}
\end{figure}

The dominant failure mode was indexing.  Twiddle factors were
repeatedly computed using the SMEM tile size instead of the global
problem size $N$, producing kernels whose tile-local FFT was
correct but whose composition into a global FFT was not. The
\texttt{debug\_correctness} tool consistently localized the
discrepancy to ``first wrong element at index 0'', which the
agent interpreted as a missing initial radix-2 stage---and added
one---without resolving the underlying indexing bug.

\subsection{Implication: Skills Decompose into Strategy and
Implementation}

We therefore organize $\mathcal{K}$ as a two-level store:

\begin{itemize}
    \item \textbf{Strategy} $\mathcal{S}^{\text{strat}}$:
    transferable, learnable from trajectories,
    LLM-discoverable. Examples: \texttt{smem\_stage\_fusion},
    \texttt{higher\_radix}, \texttt{online\_streaming\_reduce},
    \texttt{temporal\_tiling}, \texttt{fused\_scatter\_gather}.
    \item \textbf{Implementation} $\mathcal{S}^{\text{impl}}$:
    operator-specific verified code with critical-correctness
    notes (e.g., ``twiddle uses global $N$''; ``Cooley--Tukey
    in-place, not Stockham''). Provided by domain experts or
    distilled from existing libraries.
\end{itemize}

This structure refines the three claims of
Section~\ref{sec:keo} (\S\ref{sec:keo}): \textbf{Discovery}
applies at the strategy layer; \textbf{Transfer} crosses operators
at the strategy layer but not the implementation layer
(SMEM-stage fusion transfers from FFT to Stencil's temporal
tiling, but butterfly indexing does not); and \textbf{Compactness}
holds with a multiplicative twist---few strategies $\times$ many
implementations---rather than a single flat library. In our
current evaluation, 5 strategies span 10 operators while requiring
8 distinct implementations.

% ============================================================================
\section{Related Work}
% ============================================================================

\paragraph{LLM-based GPU kernel generation.}
KernelBench~\cite{ouyang2025kernelbench} benchmarks LLMs on
generating PyTorch-equivalent CUDA kernels and reports that even
state-of-the-art models produce correct-and-fast kernels on only
$\sim\!20\%$ of tasks; The AI CUDA Engineer~\cite{lange2025aicuda}
shows that a full agentic loop (translate, optimize, compose) can
substantially raise this rate but relies on a curated ``innovation
archive'' of human-written reference kernels. Other learning-to-edit
work~\cite{shypula2024kernelagent} treats kernel optimization as
sequence-to-sequence editing without a profiling environment. Our
contribution is orthogonal: rather than scaling the model or the
training data, we organize knowledge into a two-level skill store
(strategy~+~verified implementation) and ground the agent in NCU-based
profiling, allowing the same LLM to recover from incorrect
implementation attempts (e.g.\ FFT, \S\ref{sec:fft-case}) by
retrieving a verified template instead of regenerating one.

\paragraph{Compilers and autotuning.}
Halide~\cite{ragan2013halide}, TVM~\cite{chen2018tvm}, and
Ansor~\cite{zheng2020ansor} expose a structured schedule space and
search it with cost models or autotuners; Triton~\cite{tillet2019triton}
provides a tile-level programming model that our agent's codegen
backend reuses. These systems treat each kernel independently and do
not transfer optimization decisions across operators. Our skill
abstraction sits one level above schedules: a skill identifies
\emph{which} kernel-level transformation to attempt, after which a
schedule-level search such as Triton autotune or Ansor can refine the
configuration.

\paragraph{IO-aware kernel design.}
FlashAttention~\cite{dao2022flashattention} popularized the IO-aware
philosophy of fusing memory-bound passes via on-chip online algorithms.
Subsequent work---most notably Flash-KMeans~\cite{flashkmeans2026},
whose IO analysis directly inspired our Flash-GMM derivation---shows
that the same recipe applies to k-means assignment and update phases.
We generalize the recipe across the 7 Berkeley Dwarfs and quantify
\emph{when it fails}: Tab.~\ref{tab:io-vs-speedup} demonstrates that
IO reduction is necessary but not sufficient, with five
distinct failure modes (compute-bound, L2-resident data, broken GEMM
parallelism, library ceiling, atomic contention). To our knowledge,
this is the first systematic mapping of IO-aware optimization onto the
Berkeley Dwarf taxonomy.

\paragraph{Agent architectures and skill libraries.}
ReAct~\cite{yao2023react} formalized the
thought$\to$action$\to$observation loop that our agent instantiates
over the KEO interfaces. Voyager~\cite{wang2024voyager} demonstrated
that an LLM agent in Minecraft can grow a reusable skill library by
verifying behaviors in a live environment; we adapt this idea to GPU
kernels, where the verifying environment is NCU plus a correctness
oracle and the skills are GPU optimization strategies. The key novel
finding here is the \emph{two-level} structure: in Voyager, code-level
skills are themselves transferable, whereas in GPU optimization the
implementation layer (e.g.\ butterfly indexing) is operator-specific
and resists LLM regeneration even when the strategy layer is
correctly chosen.

\paragraph{Hardware profiling for optimization.}
NVIDIA Nsight Compute~\cite{nvidia2024ncu} provides per-kernel
counters (DRAM throughput, SM occupancy, warp stalls, Tensor Core
activity) that have traditionally been consumed by human experts. We
package NCU as a first-class \texttt{Environment} interface for the
agent. Tab.~\ref{tab:ncu} shows that runtime signals such as
\texttt{warps=99\%, DRAM=0\%} are diagnostic in a way that no static
analysis can reproduce, and Tab.~\ref{tab:ncu-impact} quantifies the
resulting reduction in iteration count.

% ============================================================================
\section{Cost, Limitations, and Reproducibility}
\label{sec:limits}
% ============================================================================

\subsection{Agent Cost}

Each operator costs the agent a small, bounded budget of LLM
calls, NCU profiles, and wall-clock time. Table~\ref{tab:cost}
reports per-operator means measured across the 10 evaluation
operators on a single A100 node.

\begin{table}[h]
\centering
\caption{Per-operator agent cost (mean over 10 operators).
\emph{Steps}: agent action invocations.
\emph{Tokens}: input+output tokens consumed by the LLM policy.
\emph{NCU}: number of \texttt{ncu\_profile} calls.
\emph{Wall}: end-to-end wall-clock per operator.}
\label{tab:cost}
\setlength{\tabcolsep}{4pt}
\small
\begin{tabular}{lcccc}
\toprule
Class                 & Steps & Tokens (k) & NCU & Wall (min) \\
\midrule
Memory-bound (CE, LN) & 6--9   & 18--32 & 1--2 & 3--5 \\
Algorithmic fix (KMeans) & 14--18 & 55--75 & 3--5 & 9--14 \\
End-to-end (Flash-GMM) & 28--34 & 120--160 & 6--8 & 22--30 \\
Negative (Conv2D, SpMV) & 8--11 & 25--40 & 2--3 & 4--7 \\
\bottomrule
\end{tabular}
\end{table}

For comparison, a human expert typically spends hours to days on
the same kernels; the agent's value is not in beating an expert in
isolation but in operating at scale across many kernels and
retaining the resulting skills.

\subsection{Limitations and Threats to Validity}

\paragraph{Single hardware platform.}
All speedup numbers are measured on a single NVIDIA A100 80GB
PCIe node. While the KEO interfaces are hardware-agnostic (NCU
counters, library ceilings, version tracking all generalize to
H100/H200), the \emph{numerical} thresholds in skill conditions
(L2 capacity, balance point $I^\star$, Tensor Core data types)
were tuned to A100. We expect skill \emph{rankings} to transfer
but absolute speedups to shift on H100 (larger L2, FP8 TC) and on
consumer GPUs (smaller L2, no NVLink).

\paragraph{Single LLM backend.}
The agent uses GPT-5.2 via a Copilot proxy for all experiments.
Two-level skill discovery may be sensitive to the model's prior
HPC knowledge; weaker models could fail at strategy selection
(not only implementation), collapsing the two-level distinction
that underpins our claims. We did not ablate across LLM
backends.

\paragraph{Operator sample size.}
Ten operators across seven Berkeley Dwarfs is sufficient to refute
the strong form of the IO-only hypothesis but not large enough to
estimate skill-coverage saturation. The Pareto curve in
Fig.~\ref{fig:skills-pareto} flattens after five skills on this
sample but may continue to grow on a larger benchmark.

\paragraph{Implementation library is human-curated.}
The implementation layer $\mathcal{S}^{\text{impl}}$ contains
verified templates that we wrote or distilled from cuFFT/cuBLAS.
The agent does \emph{not} discover implementations autonomously
(this is precisely the negative finding of \S\ref{sec:fft-case}).
A self-bootstrapping implementation layer remains future work.

\paragraph{Correctness oracle.}
We rely on a numerical-equivalence oracle (max abs/rel error
against a reference implementation). Kernels that are correct in
finite precision but diverge in long iterative loops (e.g.,
training-time GMM) require additional convergence checks beyond
single-kernel verification.

\subsection{Reproducibility}

Our complete code, agent traces, NCU outputs, and benchmark
harness are released at \url{https://github.com/HaibinLai/flashGMM}.
The repository contains: (i)~the KEO agent loop and tool
implementations (\texttt{io\_env/}); (ii)~the Triton/CUDA
kernels used as templates and as evaluation targets
(\texttt{flash\_gmm/}); (iii)~per-operator skill library
metadata; (iv)~scripts to regenerate every table and figure in
this paper, including the cost-vs-speedup figure
(Fig.~\ref{fig:cost}). All experiments use fixed random seeds and
report mean over 5 runs after warmup. Hardware: NVIDIA A100 80GB
PCIe, CUDA 12.4, PyTorch 2.4, Triton 3.0.

\begin{figure}[h]
\centering
\includegraphics[width=0.92\linewidth]{figures/agent_cost.pdf}
\Description{Scatter plot of per-operator wall-clock cost in minutes versus achieved speedup; memory-bound operators cluster at low cost and high speedup, algorithmic-fix operators sit at moderate cost and moderate speedup, and end-to-end Flash-GMM is the highest-cost point with above-1x speedup, while negative cases sit at low cost and below-1x speedup.}
\caption{Agent cost vs.\ achieved speedup.
Memory-bound operators are cheap and decisive; algorithmic
fixes (KMeans) cost more but pay off; negative cases
(Conv2D, SpMV) are recognized cheaply and abandoned.}
\label{fig:cost}
\end{figure}

% ============================================================================
\section{Conclusion}
% ============================================================================

We presented a framework for automatically discovering, abstracting, and transferring GPU kernel optimization skills. Our key finding is that HPC optimization knowledge has a \emph{two-level structure}: a compact set of 5--10 strategies (SMEM fusion, online reduction, higher radix, temporal tiling, fused scatter-gather) covers all 7 Berkeley Dwarfs and is reliably learnable by LLM agents from optimization trajectories, while a longer tail of operator-specific \emph{implementations} (e.g., Cooley--Tukey butterfly indexing, bit-reverse permutation) must be supplied as verified code templates. The FFT case study made this distinction concrete: across 9 rewrites the agent chose the right strategies but never produced correct indexing. By grounding the optimization agent in hardware profiling (NCU) and structuring its reasoning through the Knowledge--Environment--Objective framework, we achieve reliable speedups of 1.2--6.0$\times$ while correctly avoiding unprofitable optimizations. The Flash-GMM case study confirms that the same strategy library drives both isolated-kernel optimization and a complete end-to-end application, with NCU-guided refinement contributing an additional $\sim\!10\times$ kernel-internal speedup atop the algorithmic gains. The two-level skill organization---few strategies $\times$ many verified implementations---transforms kernel optimization from a per-operator search problem into a growing knowledge system that improves with experience.

\bibliographystyle{ACM-Reference-Format}
\bibliography{reference}

\end{document}
